%--- iliad ---
% title: Condensation
% summary: >-
%   Organizing information across a family of observations. Define latent
%   models, examine reconstruction and conditional independence, and prove
%   an exact and quantitative correspondence theorem. A worked
%   hierarchical-clustering example, and how condensation relates to earlier
%   work in statistics and what remains between the theorem and useful
%   translation between the latents of human and AI models.
%--- end ---
\documentclass[11pt]{article}
\IfFileExists{iliad.sty}{\usepackage[boxes]{iliad}}{\usepackage[boxes]{../iliad}}
\usepackage[round]{natbib}
\newcommand{\PI}{\mathcal{P}^{+}(I)}
\newcommand{\Up}{\mathord{\supseteq}}
\title{Condensation}
\author{Satya Benson}
\date{21 September 2026}
\begin{document}
\maketitle

\youtube{odsHleu80x8}

\section{Prerequisites}
\label{sec:prerequisites}
\begin{enumerate}

\item We use probability on finite sets, conditional probability, and elementary
proofs.
\item All random variables in this sheet have finite ranges.
\item All logarithms
have base two.

% \item We use the latent-model definition and conditional
% entropy from the main worksheet.
%The example adapts the hierarchy in Jeremy Gillen and Daniel
%Chiang's \emph{A summary of Condensation and its relation to Natural Latents},
%\S1.1.1. The construction and figure here are new; we do not use their numerical
%entropy estimates.

%This reading follows the condensation worksheet.
\item We use conditional entropy,
conditional independence, and the distinction between an individual latent and
a collection of latents.
%All variables in the examples have finite alphabets;
%logarithms are base two.
\end{enumerate}


\begin{learningoutcomes}
\subsection*{Condensation}
\begin{itemize}
\item Construct a latent model and distinguish LVM from reconstruction.
\item Explain what the Markov condition excludes, using a concrete example.
\item Prove correspondence between collections of latents, first exactly and
then with explicit error terms.
\item Distinguish the correspondence theorem from claims about learning,
computationally finding a translation, or choosing observables.
\end{itemize}
\subsection*{A hierarchical clustering example}
\begin{itemize}
\item Distinguish LVM, where the latents assigned to an observation determine
it, from reconstruction, where an observation recovers those latents.
\item Compare the information charged by different latent models for the same query.
\end{itemize}
\subsection*{Condensation in context}
\begin{itemize}
\item Relate condensation to sufficient statistics, common information,
computational mechanics, and graphical models.
\item Separate informational correspondence from finding and using a translation.
\item Identify which additional claims would connect correspondence to AI safety.
\end{itemize}
\end{learningoutcomes}

\subsection*{How to use this sheet}
Read \cref{sec:questions,sec:conditions} after the first lecture, then work on
\crefrange{ex:bits}{ex:duplicates}. After the second lecture, read
\cref{sec:correspondence} and complete \crefrange{ex:witness}{ex:induction}. These eight exercises form the
core route. \Cref{ex:blocks,ex:scores} are optional extensions.

\subsection{Information theory reference}
\label{sec:it-reference}
For a finite-valued variable $U$,
\[
 H(U)=-\sum_u p(u)\log p(u),\qquad
 H(U\mid V)=\sum_v p(v)H(U\mid V=v),
\]
with $0\log 0=0$. Conditional entropy is nonnegative. It is zero precisely when
$U$ is a function of $V$ almost surely: values on probability-zero events do not
matter. The chain rule and conditioning inequality are
\[
 H(U,V\mid W)=H(U\mid W)+H(V\mid U,W),\qquad
 H(U\mid V,W)\le H(U\mid W).
\]
Conditional mutual information is
\[
 I(U;V\mid W)=H(U\mid W)-H(U\mid V,W)\ge0.
\]
It is symmetric in $U,V$ and is zero precisely when $U,V$ are conditionally
independent given $W$. Omitting $W$ gives mutual information $I(U;V)$.
Entropy of a tuple means joint entropy, not the sum of its marginal entropies.
An empty tuple is constant and has zero entropy.
For probability distributions $p,q$ on a finite set,
\[
 D_{\mathrm{KL}}(p\Vert q)=\sum_u p(u)\log\frac{p(u)}{q(u)}\ge0.
\]
This last quantity appears only in the optional score discussion.

\section{Condensation}
\subsection{A family of questions}
\label{sec:questions}
Fix jointly distributed observables $X_1,\ldots,X_n$. We look for latent
variables, each assigned to a set of observations, such that the latents
assigned to an observation determine it.
We will give conditions under which corresponding latent collections in two
such models determine one another, and bound the remaining conditional entropy
when the conditions hold approximately. We take the joint distributions as given.

\begin{remark}[Choosing the observables]
One can write $X_i=f_i(S)$ for an underlying variable $S$ representing sense
data or a physical state. This introduces no restriction: $S$ could simply be
the tuple of all observations. An observable distinguishes outcomes by
partitioning them into its possible answers; random-variable values label those
parts. Information quantities do not depend on bijective relabeling of answers.
The choice of which observables to include, and their joint law, remains part
of the model.
\end{remark}

\begin{example}[Overlapping observations]
\label{eg:bits}
Let $S,P,Q$ be independent fair bits and set
\[
 X_1=(S,P),\qquad X_2=(S,Q),\qquad X_3=(P,Q).
\]
Use a latent $Y_{\{1,2\}}=S$ for observations 1 and 2, a latent
$Y_{\{1,3\}}=P$ for observations 1 and 3, and a latent $Y_{\{2,3\}}=Q$
for observations 2 and 3. Each observation is determined by the two latents
assigned to it. Each of these latents can also be recovered from either
observation to which it is assigned.
\end{example}

\subsubsection{Addresses and collections}
Let $I=\{1,\ldots,n\}$ and write $\PI$ for its nonempty subsets.
An \emph{address} $A\in\PI$ specifies which observations receive a latent
$Y_A$. The collection of addresses is a partially ordered set under inclusion.
We use one slot for each nonempty subset, setting unused slots to constants.

\begin{definition}[Latent model]
\label{def:model}
A latent model of $X=(X_i)_{i\in I}$ consists of finite-valued variables
$Y=(Y_A)_{A\in\PI}$ jointly distributed with $X$, such that
\begin{equation}
 H\bigl(X_i\mid Y_{\{A\in\PI\,:\,i\in A\}}\bigr)=0
 \quad\text{for every }i\in I.
 \label{eq:forward}
\end{equation}
Here $Y_{\mathcal{F}}=(Y_A)_{A\in\mathcal{F}}$ for a family
$\mathcal{F}\subseteq\PI$.
\end{definition}

We call \eqref{eq:forward} the \emph{latent variable model condition}, or
\emph{LVM}. The contributing latents, taken together, determine the
observation. This is not a claim that an agent knows their realized values. Some slots may represent residual noise.
Also, an address need not be minimal: being assigned to an observation does not
assert that a slot makes a nonredundant or causal contribution to it.

For nonempty $A\subseteq I$ and $B\subseteq I$, define two different families:
\[
 \Up A=\{C\in\PI:A\subseteq C\},\qquad
 \mathcal{J}_B=\{C\in\PI:C\cap B\ne\varnothing\}.
\]
The tuple $Y_{\Up A}$ is the \emph{tower above $A$}. The tuple
$Y_{\mathcal{J}_B}$ contains all latents contributing to at least one observation
in $B$, and therefore determines $X_B=(X_i)_{i\in B}$.
In particular $\Up\{i\}=\mathcal{J}_{\{i\}}$ is the family of addresses
contributing to $X_i$, so \eqref{eq:forward} reads $H(X_i\mid Y_{\Up\{i\}})=0$.
For a larger $A$, $\Up A$ and $\mathcal{J}_A$ generally differ.

\begin{exercise}[Overlapping bits; 15 minutes]
\label{ex:bits}\important
Use \cref{eg:bits}, with all unmentioned slots constant.
\begin{enumerate}
\item Compute $H(X_1)$, $H(X_1,X_2,X_3)$, and $I(X_1;X_2)$.
\item List $\Up\{1\}$, $\Up\{2\}$, and $\Up\{1,2\}$, including
constant slots. Which tuple determines $X_1$? Which of the three is
$\Up\{1\}\cap\Up\{2\}$?
\item Give explicit functions recovering $X_1$ from its contributing latents
and recovering $Y_{\{1,2\}}$ from $X_1$.
\item Does the address structure have to be a tree? Identify two overlapping
addresses in this example neither of which contains the other.
\end{enumerate}
\end{exercise}

\subsection{Two conditions on a latent model}
\label{sec:conditions}
LVM is not one of the two conditions below. It is part of the definition of a
latent model, and every observable law admits one by construction: put a copy
of each $X_i$ in its singleton slot, or the entire tuple $X$ in the top slot
$Y_I$, and LVM holds either way. Satisfying it therefore says nothing about how
information is shared. The two conditions below constrain these alternatives in
different ways.

\begin{definition}[Reconstruction]
\label{def:reconstruction}
A latent model has exact reconstruction if
\[
 H(Y_A\mid X_i)=0\qquad\text{whenever }i\in A.
\]
Equivalently, $H(Y_{\Up A}\mid X_i)=0$ whenever $i\in A$: every slot
in that finite tuple can be recovered from $X_i$.
\end{definition}

The equivalence in this definition uses exact zero errors. In approximate
statements we will bound the entropy of the \emph{whole tower}; separately small
errors for its individual slots could add up.

A family $\mathcal{F}\subseteq\PI$ is \emph{upward closed} if
$A\in\mathcal{F}$ and $A\subseteq C$ imply $C\in\mathcal{F}$.
For example, $\Up A$ and $\mathcal{J}_B$ are upward closed. Taking a
collection together with every slot above it prevents us from omitting higher
variables on which it may depend.

\begin{definition}[Ordered Markov condition and defect]
\label{def:markov}
The ordered Markov condition is
\[
 I(Y_{\mathcal{F}};Y_{\mathcal{G}}\mid
 Y_{\mathcal{F}\cap\mathcal{G}})=0
\]
for all upward-closed families $\mathcal{F},\mathcal{G}\subseteq\PI$.
Its quantitative defect is
\[
 \eta_Y=\max_{\mathcal{F},\mathcal{G}\text{ upward closed}}
 I(Y_{\mathcal{F}};Y_{\mathcal{G}}\mid Y_{\mathcal{F}\cap\mathcal{G}}).
\]
\end{definition}

The conditioning family $\mathcal{F}\cap\mathcal{G}$ is an intersection of
\emph{address families}, not a quantity inferred to be common by inspecting
their values. The condition says that the latents at those shared addresses
account for the dependence between the two collections. It does not require
every latent to be independent of every other latent. In the two-observable
case, its only potentially nonzero requirement is
\[
 I(Y_{\{1\}};Y_{\{2\}}\mid Y_{\{1,2\}})=0.
\]
This is a graphical-model condition for the inclusion order. The equivalent
local version says that each slot is conditionally independent of all
incomparable slots given its strict supersets.

\begin{definition}[Perfect condensation]
\label{def:perfect}
A latent model is a perfect condensation if it has exact reconstruction
and satisfies the ordered Markov condition.
\end{definition}
This is equivalent to Eisenstat's definition by optimal conditioned scores
\citep[Theorem~5.10]{eisenstat2025condensation}. \Cref{sec:scores} explains that connection.
The overlapping-bit model in \cref{eg:bits} is a perfect condensation.

\begin{exercise}[Two failed models; 20 minutes]
\label{ex:failed}\important
Keep the observables of \cref{eg:bits}. Consider two alternative models.
\begin{enumerate}
\item Set $Y_{\{i\}}=X_i$ for $i=1,2,3$ and every other slot constant.
Check LVM and reconstruction. Take
$\mathcal{F}=\Up\{1\}$ and $\mathcal{G}=\Up\{2\}$.
Calculate the conditional mutual information in \cref{def:markov}.
Which information is duplicated across singleton slots?
\item Instead set $Y_{\{1,2,3\}}=(S,P,Q)$ and every other slot constant.
Check LVM and the Markov condition. Compute
$H(Y_{\{1,2,3\}}\mid X_1)$. Which condition fails?
\item Explain why LVM, reconstruction, and the
Markov condition have different roles.
\end{enumerate}
\end{exercise}

\subsubsection{Existence is a separate question}
\label{sec:existence}
The preceding examples test particular latent models. A stronger question is
whether a given observable law admits \emph{any} perfect condensation.
Even a pair of correlated binary variables can fail to do so.

Let $I=\{1,2\}$, let $X_1$ be a fair bit, and set $X_2=X_1\mathbin{\oplus}N$,
where $N$ is independent Bernoulli$(q)$ noise and $0<q<1/2$. Here
$\oplus$ is addition modulo two, so $a\oplus b$ is $0$ when the two bits agree
and $1$ when they differ.
The joint law is
\[
\renewcommand{\arraystretch}{1.8}
\begin{array}{r|cc}
 P(x_1,x_2) & \;\;x_2=0\;\; & \;\;x_2=1\;\;\\\hline
 x_1=0 & \dfrac{1-q}{2} & \dfrac{q}{2}\\
 x_1=1 & \dfrac{q}{2} & \dfrac{1-q}{2}
\end{array}
\]
Every one of the four outcomes has positive probability, because $0<q<1/2$;
and $X_1,X_2$ are dependent, because the two rows differ.
Suppose a perfect condensation existed, and consider its top slot
$Y_{\{1,2\}}$.

With two observations there are only three addresses, so the whole model is
three slots:
\[
\begin{array}{r|ccc}
 \text{address} & \{1\} & \{1,2\} & \{2\}\\
 \text{slot} & Y_{\{1\}} & Y_{\{1,2\}} & Y_{\{2\}}
\end{array}
\]
LVM says $X_1$ is determined by $(Y_{\{1\}},Y_{\{1,2\}})$ and $X_2$ by
$(Y_{\{2\}},Y_{\{1,2\}})$. The top slot $Y_{\{1,2\}}$ is the only one assigned to both
observations, so it is the only place shared information can sit.

Reconstruction would give $Y_{\{1,2\}}=f(X_1)=g(X_2)$ almost surely. Because
all four observable pairs have positive probability, these equalities force
\[
 f(0)=g(0)=g(1)=f(1).
\]
Thus $Y_{\{1,2\}}$ is constant. The Markov condition would then require
$Y_{\{1\}}$ and $Y_{\{2\}}$ to be independent. LVM
makes each $X_i$ a function of its singleton slot and the constant $Y_{\{1,2\}}$.
Functions of independent variables are independent, contradicting the chosen
observable law. No perfect condensation exists for this pair.

There is a quantitative obstruction if we retain \emph{exact} reconstruction. The same argument makes $Y_{\{1,2\}}$ constant. Data processing then
gives
\[
 \eta_Y=I(Y_{\{1\}};Y_{\{2\}})
 \ge I(X_1;X_2)=1-h_2(q)>0,
\]
where $h_2(q)=-q\log q-(1-q)\log(1-q)$ is the binary entropy function.
Dependence between the observables must remain as a Markov defect when their
separately recoverable top slot is constant.

That bound pins one end of a trade-off: exact reconstruction forces a Markov
defect of at least $1-h_2(q)$, and tolerating some reconstruction error buys
part of it back. When no perfect condensation exists, which is the usual
situation, the latent models of a given observable law trace out a Pareto
frontier between reconstruction defect and Markov defect. The useful question
is then whether some point on that frontier has both low \emph{enough} for the
correspondence bound of \cref{sec:correspondence} to say something.

\begin{remark}
One could instead ask for the model minimizing the conditioned score of
\cref{sec:scores}. The score is defined per observation block, however, and
there is not always a single model minimizing it everywhere.
\end{remark}

\begin{exercise}[Collections, not individual slots; 15 minutes]
\label{ex:duplicates}\important
For this exercise let $I=\{1,2\}$, take independent fair bits $C,R,T$, and set
$X_1=(C,R)$, $X_2=(C,T)$. Compare
\[
\begin{array}{lll}
 Y_{\{1,2\}}=C,&Y_{\{1\}}=R,&Y_{\{2\}}=T,\\
 Z_{\{1,2\}}=C,&Z_{\{1\}}=(C,R),&Z_{\{2\}}=(C,T).
\end{array}
\]
\begin{enumerate}
\item Check that both models are perfect condensations.
\item Compute $H(Z_{\{1\}}\mid Y_{\{1\}})$.
\item Show that $Y_{\Up\{1\}}$ and $Z_{\Up\{1\}}$ determine one
another. Explain why correspondence between these tuples is a different
claim from correspondence between the singleton slots.
\end{enumerate}
\end{exercise}

\subsection{The correspondence theorem}
\label{sec:correspondence}
Take two latent models of the same observable tuple $X$. To compare them,
put $(X,Y,Z)$ on one probability space while retaining both specified laws
$(X,Y)$ and $(X,Z)$. For finite variables this is always possible: for example,
use $p(x,y,z)=p(x)p(y\mid x)p(z\mid x)$. The theorem holds for any such
coupling.

\begin{theorem}[Correspondence]
\label{thm:correspondence}
For two latent models $Y,Z$ of $X$ and any nonempty $A\subseteq I$, let
$k=|A|$. Then
\begin{equation}
 H(Y_{\Up A}\mid Z_{\Up A})
 \le \sum_{i\in A}H(Y_{\Up A}\mid X_i)+(k-1)\eta_Z.
 \label{eq:correspondence}
\end{equation}
If each term in the sum is at most $\delta$ and $\eta_Z\le\eta$, the bound
is $k\delta+(k-1)\eta$. If both budgets are $\varepsilon$, it is
$(2k-1)\varepsilon$. Interchanging $Y,Z$ gives the reverse bound.
\end{theorem}

\begin{corollary}[Exact correspondence]
\label{cor:exact}
If $Y$ and $Z$ are perfect condensations, their towers above every nonempty
address $A$ determine one another almost surely:
\[
 H(Y_{\Up A}\mid Z_{\Up A})=H(Z_{\Up A}\mid Y_{\Up A})=0.
\]
\end{corollary}

The forward bound uses reconstruction in the sender $Y$, LVM in the receiver
$Z$, and the Markov defect of $Z$.
It does not use the Markov condition of $Y$ until we reverse the argument.
The result concerns corresponding collections. \Cref{ex:duplicates} shows why we cannot simply
replace those collections by individual slots.

This proof is a specialization of Eisenstat's approximate correspondence
theorem \citep[Theorem~6.8]{eisenstat2025condensation}, using the two-witness
inequality \citep[Lemma~6.4]{eisenstat2025condensation}.
Its tower presentation follows the exposition listed in \cref{sec:reading}.
The exercises below give the full argument without additional machinery.

\subsubsection{Proof roadmap}
The theorem has three steps.
\begin{enumerate}
\item If two witnesses each determine a target and are independent given their
common part, that common part determines the target. Entropy gives a version
with errors.
\item Each observation $X_i$ is determined by its contributing $Z$-latents.
Anything recoverable from $X_i$ is therefore recoverable from that collection.
\item Intersect these contributing collections over $i\in A$. Exactly the
addresses containing all of $A$ survive. Each intersection costs at most one
Markov error term.
\end{enumerate}
The word ``witness'' below means a variable or tuple from which the target is
recoverable.

\begin{lemma}[Two witnesses]
\label{lem:witness}
For any finite-valued variables $U,M,P,Q$,
\begin{equation}
 H(U\mid M)\le H(U\mid M,P)+H(U\mid M,Q)+I(P;Q\mid M).
 \label{eq:witness}
\end{equation}
\end{lemma}

\begin{exercise}[Prove the two-witness lemma; 20 minutes]
\label{ex:witness}\important
\begin{enumerate}
\item Prove \cref{lem:witness} by expanding conditional entropies using the
chain rule. Find an expression for the difference between the right and left
sides that is a sum of nonnegative quantities.
\item Deduce the exact statement when the three terms on the right are zero.
\item Explain the three error terms. Which measure failures of reconstruction,
and which measures dependence that remains after revealing $M$?
\end{enumerate}
\end{exercise}
\begin{hint}
Try expanding $I(P;Q\mid M)-I(P;Q\mid U,M)$, then rearranging with the
chain rule, which is stated in
\cref{sec:it-reference}.
The desired difference contains a conditional mutual information
and a conditional entropy.
\end{hint}

\begin{exercise}[The dependence term is needed; 5 minutes]
\label{ex:dependence}
Let $U=P=Q$ be a fair bit and let $M$ be constant. Evaluate every term in
\eqref{eq:witness}. What would go wrong if the mutual-information term were
omitted? How does this example resemble the singleton-copy model in \cref{ex:failed}?
\end{exercise}

\begin{exercise}[Transfer through an observable; 10 minutes]
\label{ex:transfer}\important
Let $U$ be any finite-valued target and suppose $X_i$ is determined by
$Z_{\Up\{i\}}$. Prove
\[
 H(U\mid Z_{\Up\{i\}})\le H(U\mid X_i).
\]
State the equality obtained by adding $X_i$ to the conditioning tuple, then
apply the conditioning inequality. Does this argument assume that $U$ and
$Z_{\Up\{i\}}$ are conditionally independent given $X_i$?
\end{exercise}

\begin{exercise}[Intersect the address families; 15 minutes]
\label{ex:intersections}\important
\begin{enumerate}
\item For $I=\{1,2,3\}$, list $\Up\{1\}$, $\Up\{2\}$, and
$\Up\{3\}$. Compute $\Up\{1\}\cap\Up\{2\}$ and then its
intersection with $\Up\{3\}$.
\item Show that an intersection of upward-closed families is upward closed.
\item For general nonempty $A\subseteq I$, prove
$\bigcap_{i\in A}\Up\{i\}=\Up A$ by checking membership of an address
$C$ on both sides.
\item If $\mathcal{F},\mathcal{G}$ are upward closed, apply the two-witness
lemma with $M=Z_{\mathcal{F}\cap\mathcal{G}}$,
$P=Z_{\mathcal{F}}$, and $Q=Z_{\mathcal{G}}$.
Why can the extra conditioning on $M$ be dropped from each reconstruction
term? Write the resulting bound using $\eta_Z$.
\end{enumerate}
\end{exercise}

\begin{exercise}[Complete the correspondence proof; 25 minutes]
\label{ex:induction}\important
Fix $A=\{i_1,\ldots,i_k\}$, set
$\mathcal{G}_j=\bigcap_{\ell=1}^{j}\Up\{i_\ell\}$.
\begin{enumerate}
\item Use \cref{ex:transfer} to prove the bound at $j=1$.
\item Use \cref{ex:intersections} to show, for $j\ge2$,
\[
 H(Y_{\Up A}\mid Z_{\mathcal{G}_j})\le
 H(Y_{\Up A}\mid Z_{\mathcal{G}_{j-1}})+H(Y_{\Up A}\mid Z_{\Up\{i_j\}})+\eta_Z.
\]
\item Prove by induction that
\[
 H(Y_{\Up A}\mid Z_{\mathcal{G}_j})\le
 \sum_{\ell=1}^{j}H(Y_{\Up A}\mid X_{i_\ell})+(j-1)\eta_Z.
\]
Evaluate this at $j=k$ to prove \cref{thm:correspondence}. Check the
case $k=1$ explicitly.
\item Deduce \cref{cor:exact}, including the reverse direction. Name the
assumption used at each step of the proof.
\item Suppose instead that $H(Y_C\mid X_i)\le\varepsilon$ separately for
all relevant slots $C$. Why does this not by itself supply the tower error
budget $H(Y_{\Up A}\mid X_i)\le\varepsilon$?
\end{enumerate}
\end{exercise}

\subsection{Optional: several observations as one witness}
\label{sec:blocks}
A shared parameter need not be recoverable from a single observation. For
example, take $n\ge2$, let $C$ be a fair bit, and set
$X_i=C\mathbin{\oplus}N_i$, where the
$N_i$ are independent Bernoulli$(p)$ noises with $0<p<1/2$, independent of
$C$. The model $Y_I=C$, $Y_{\{i\}}=X_i$ satisfies LVM and the
Markov condition exactly, but $H(C\mid X_i)=h_2(p)>0$, where
$h_2$ is the binary entropy function of \cref{sec:existence}.
Larger observation blocks can reveal more about $C$. To use them in the proof,
replace $\Up\{i\}$ by $\mathcal{J}_B$.

\begin{exercise}[Block witnesses; optional, 25 minutes]
\label{ex:blocks}
\begin{enumerate}
\item Let $I=\{1,2,3\}$ and suppose a target $U$ has conditional entropy at
most $\delta$ given each of $X_{\{1,2\}}$, $X_{\{1,3\}}$, and
$X_{\{2,3\}}$. Which addresses survive the intersection
$\mathcal{J}_{\{1,2\}}\cap\mathcal{J}_{\{1,3\}}\cap\mathcal{J}_{\{2,3\}}$?
Prove a bound on the entropy of $U$ given the surviving $Z$-slots.
\item More generally, let $\mathcal{W}$ be a nonempty finite family of nonempty
observation blocks. Adapt \cref{ex:induction} to show
\[
 H(U\mid Z_{\mathcal{R}})\le
 \sum_{B\in\mathcal{W}}H(U\mid X_B)+(|\mathcal{W}|-1)\eta_Z,
 \qquad \mathcal{R}=\bigcap_{B\in\mathcal{W}}\mathcal{J}_B.
\]
\item Fix an integer $0\le r<|A|$ and take every $(r+1)$-element subset of $A$ as a
witness. Show that $C\in\mathcal{R}$ exactly when $|A\setminus C|\le r$.
What does this give for $r=0$? Why is the receiver generally larger for $r>0$?
\end{enumerate}
\end{exercise}

Eisenstat's later almost-perfect-condensation formulation
\citep{eisenstat2025condensation} bounds the tower
reconstruction error whenever $|A\cap B|>r$, the Markov defect, and the number
of active slots. The threshold allows small slots to account for local noise
without requiring every such slot to be recovered from every relevant block.
It also changes which receiver the proof isolates, as \cref{ex:blocks} shows. The exact-tower
bound with singleton witnesses does not follow from an arbitrary positive
threshold $r$.

\subsection{Optional: the conditioned score}
\label{sec:scores}
The conditions also have a description-length interpretation. For an
observation block $B$, compare
\[
 \sigma(B)=\sum_{A\in\mathcal{J}_B}H(Y_A),\qquad
 \chi(B)=\sum_{A\in\mathcal{J}_B}H(Y_A\mid Y_{\Up A\setminus\{A\}}).
\]
The simple score $\sigma$ charges each active slot separately. The conditioned
score $\chi$ charges its additional information given strict supersets.
Those supersets are also active for $B$. Encoding in reverse inclusion order
makes this conditioning available. These are ideal expected description-length
quantities; they do not measure computational retrieval time.

Both scores are at least $H(X_B)$. More precisely, put
$V_B=\chi(B)-H(Y_{\mathcal{J}_B})$. Then
\begin{equation}
 \chi(B)-H(X_B)=V_B+H(Y_{\mathcal{J}_B}\mid X_B),\qquad V_B\ge0.
 \label{eq:score}
\end{equation}
The second term charges latent information not recoverable from $X_B$.
The first charges dependence not accounted for by the inclusion hierarchy.
It vanishes under the ordered Markov condition.

For a graphical-model expression, define
\[
 q_B(y)=\prod_{A\in\mathcal{J}_B}
 p(y_A\mid y_{\Up A\setminus\{A\}}).
\]
This is a normalized distribution: sample larger addresses before smaller
ones. Conditional kernels on probability-zero parent configurations can be
chosen arbitrarily. Direct expansion gives
$V_B=D_{\mathrm{KL}}(p_{Y_{\mathcal{J}_B}}\Vert q_B)$.
Thus \eqref{eq:score} separates two nonnegative defects.

\begin{exercise}[The two charges; optional, 25 minutes]
\label{ex:scores}
\begin{enumerate}
\item Order $\mathcal{J}_B$ so strict supersets appear first. Use the chain rule
and conditioning inequality to prove $\chi(B)\ge H(Y_{\mathcal{J}_B})$.
\item Use LVM to prove
$H(Y_{\mathcal{J}_B})-H(X_B)=H(Y_{\mathcal{J}_B}\mid X_B)$, hence
\eqref{eq:score}.
\item For both models in \cref{ex:duplicates}, calculate $\sigma(\{1,2\})$ and
$\chi(\{1,2\})$. Why is copying the top variable into a lower slot free
under one score but not the other?
\item If $A\cap B\ne\varnothing$, show that
$H(Y_{\Up A}\mid X_B)\le\chi(B)-H(X_B)$.
\end{enumerate}
\end{exercise}

Eisenstat defines perfect condensation by $\chi(B)=H(X_B)$ for every $B$.
The equivalence with the characterization above is his Theorem 5.10.
The decomposition above also explains why small score gaps imply strong
reconstruction conditions. The converse for positive-threshold almost-perfect
condensation requires additional reconstruction coverage: some active slots
may never meet that threshold. Do not identify the two approximate criteria
without checking these quantifiers.

\subsection{What the result says, and connections}
\label{sec:meaning}
The theorem compares two joint laws extending the same observable law. Under
the exact conditions, the two latent towers at each address determine one
another almost surely; under the quantitative hypotheses, their conditional
entropies satisfy the stated bounds. The theorem does not select the observable
family or give a learning algorithm. A perfect condensation need not exist,
as the full-support example in \cref{sec:existence} shows.

\subsubsection{Connections to established questions}
In statistics, a statistic is sufficient for a parameter when it
retains the data's information relevant to that parameter. Predictive
sufficiency instead concerns retaining information about a target to be
predicted. Condensation studies the organization of information across an
indexed family of observations, with explicit contribution addresses and a
correspondence result. It shares questions about information retention with
sufficiency, but these are not identical definitions or claims.

Graphical models supply the language of conditional independence and
factorization used by the ordered Markov condition. The graph here is fixed by
inclusion of addresses; its edges do not by themselves have a causal meaning.
Common-information questions also ask what two variables share. Our
correspondence argument applies across a family of overlapping address
collections, using the same entropy calculus rather than a new information
measure. Computational mechanics studies predictive states obtained from
histories with the same conditional future law. Condensation's formulation is
static: time, learning, and a prescribed past/future split are not required.
None of these comparisons implies that earlier frameworks cannot study
structured or modular organizations of information.

\subsubsection{Limits of information-theoretic correspondence}
If $V=f(U)$ for a bijection $f$, then
$H(U\mid V)=H(V\mid U)=0$. This remains true when describing or discovering
$f$ is expensive. The theorem does not provide an efficient translation
algorithm or account for the description length of the probability law.

Small conditional entropy is an average guarantee. If $E$ is Bernoulli$(p)$,
$D$ is an independent fair bit, $U=(E,ED)$, and $V=E$, then $H(U\mid V)=p$.
On the rare event $E=1$, the missing bit is still completely unknown.
Conversely, a rare event with an enormous number of outcomes can dominate an
entropy. A safety application needs to say which failures matter, not just
report an average information budget.

\subsubsection{Discussion prompts}
\begin{enumerate}
\item Choose a comparison from
\cref{sec:context-connections}. State one mathematical similarity and
one difference in what is being characterized. If it relates to one of the
condensation conditions, say how.
\item Corresponding latent collections in two models determine one another with
zero conditional entropy. What does this establish, and what else might be
needed to find and use the translation between them?
\item Where might a correspondence like this matter for AI safety? Trace one
chain from the theorem to something you would want, and say where it breaks:
which steps does the theorem not establish?
\Cref{sec:context-safety} sets out one such chain.
\end{enumerate}

\section{A hierarchical clustering example}
\subsection{Thirty-six observations}
There are two groups, each containing two clusters, each with nine observations.
Write $X_{gcj}$ for the $j$th two-dimensional point in cluster $c$ of group $g$.
The labels $(g,c,j)$ and memberships are fixed. We are comparing models of this
labeled family, not inferring an unknown partition of an unlabeled cloud.

\begin{figure}[ht]
\centering
\includegraphics[width=\linewidth]{fig/hierarchical-points.pdf}
\caption{One draw from the finite model. The left panel shows the 36 points;
the right panel shows which set of observations receives each parameter, and is
not a graph inferred from the points. The global shape $S$ is available to every
observation.}
\label{fig:hierarchy}
\end{figure}

The parameters are a global shape $S$, group centers $G_1,G_2$, and cluster
centers $M_{gc}$. Given these parameters, the 36 observations are independent.
Marginalizing over shared parameters generally makes observations dependent.
The plotted cloud is one realization; the random variables and their entropies
refer to the joint law across realizations.

\subsection{An explicit finite law}
Let $D=\{0,\ldots,200\}^2$. Draw $S$ uniformly from three shape labels, and
draw $G_1,G_2$ independently and uniformly from $\{60,\ldots,140\}^2$,
independently of $S$. Conditional on $G_g$, draw its two centers independently:
\[
 P(M_{gc}=m\mid G_g=a)\ \propto\
 \exp\!\left(-\frac{\|m-a\|^2}{128}\right),\qquad m\in D.
\]
Given $S$ and all centers, draw the observations independently, using
\[
 P(X_{gcj}=x\mid M_{gc}=m,S=s)\ \propto\
 \exp\!\left(-\tfrac12(x-m)^\mathsf{T}Q_s^{-1}(x-m)\right),\quad x\in D,
\]
where
\[
 Q_{\mathrm{round}}=\begin{pmatrix}5&0\\0&5\end{pmatrix},\quad
 Q_{\mathrm{right}}=\begin{pmatrix}25&20\\20&25\end{pmatrix},\quad
 Q_{\mathrm{left}}=\begin{pmatrix}25&-20\\-20&25\end{pmatrix}.
\]
Each proportionality means normalization by the sum over $D$. These are
probability mass functions on a finite grid. The matrices control shape; they
are not the exact covariances after discretization and boundary truncation.
Every grid point has positive probability. Cluster centers tend to lie near
their group center, but the four components need not be visually separated in
every realization.

\subsection{A structured latent model}
Let $I$ contain all 36 indices. Let $A_g$ be the 18 indices of group $g$, and
$A_{gc}$ the nine indices of cluster $(g,c)$. Assign
\[
 Y_I=S,\qquad Y_{A_g}=G_g,\qquad Y_{A_{gc}}=M_{gc},\qquad
 Y_{\{i\}}=X_i,
\]
and make every other slot constant. The joint law factorizes in the address
inclusion order, so its ordered Markov defect is zero.

The singleton slot supplies the sampled point itself. Thus
$H\bigl(X_i\mid Y_{\{A\,:\,i\in A\}}\bigr)=0$ exactly. The cluster center and shape alone
specify a conditional distribution for the point; they do not determine its
realized coordinates. The residual conditional entropy is
$H(X_i\mid M_{gc},S)$.

Reconstruction is different. All the kernels have full support, so any
finite observation block leaves several parameter configurations with positive
posterior probability. For example, $H(M_{gc}\mid X_i)>0$. The model's joint
law satisfies the Markov condition exactly, but it is not a perfect condensation.

For a nonempty block $B$, the general score-gap identity \eqref{eq:score} is
\[
 \chi(B)-H(X_B)
 =\bigl[\chi(B)-H(Y_{\mathcal J_B})\bigr]
   +H(Y_{\mathcal J_B}\mid X_B).
\]
Here the bracketed Markov term is zero. The gap is the latent information
remaining uncertain after observing the queried block.

\subsection{Three organizations of the same observable law}
For a nonempty queried \emph{observation block} $B$, the conditioned score
sums the charges of all slots whose addresses meet $B$.

\paragraph{One global tuple.}
Put $Z_I=X_I$ and make every other slot constant. The global slot determines
each observation. For every nonempty $B$,
\[
 \chi_{\mathrm{global}}(B)=H(X_I),\qquad
 \chi_{\mathrm{global}}(B)-H(X_B)=H(X_I\mid X_B).
\]
This attains the entropy bound when querying the whole tuple. Even a singleton
query incurs the entire tuple's entropy in this score.

\paragraph{An ordered chain.}
Choose an ordering $X_1,\ldots,X_{36}$. Set
$Z_{\{j,j+1,\ldots,36\}}=X_j$ and make other slots constant. The contributing
slots for $X_i$ contain $X_1,\ldots,X_i$. The charge for slot $j$ is
$H(X_j\mid X_1,\ldots,X_{j-1})$, hence
\[
 \chi_{\mathrm{chain}}(B)=H(X_1,\ldots,X_{\max B}).
\]
This attains the entropy bound for prefix queries. Querying only $X_{36}$ costs
$H(X_I)$. These statements concern this particular address assignment, not
every autoregressive architecture.

\paragraph{The structured model.}
The score charges the shape once, the group and cluster parameters needed by
$B$, and the residual uncertainty of each requested observation:
\[
\begin{aligned}
 \chi_{\mathrm{structured}}(B)={}&H(S)
 +\sum_{g:B\cap A_g\ne\varnothing}H(G_g)\\
 &+\sum_{g,c:B\cap A_{gc}\ne\varnothing}H(M_{gc}\mid G_g)
 +\sum_{i\in B}H(X_i\mid M_{g(i)c(i)},S).
\end{aligned}
\]
This expression avoids charging for unrelated observations. It need not be
smaller than the other scores for every block. In particular, for $B=I$ it
exceeds $H(X_I)$ by the remaining parameter uncertainty, whereas both baseline
constructions attain $H(X_I)$. Fitting the observable law does not by itself
settle which organization is preferable across queries.

\subsection{Why blocks help recover parameters}
For nonempty $B$ inside a fixed cluster $(g,c)$, set $\Theta=(S,G_g,M_{gc})$. The
nonconstant contributing tuple is $(\Theta,X_B)$, so
\[
 \chi_{\mathrm{structured}}(B)-H(X_B)=H(\Theta\mid X_B).
\]
For nonempty $B\subseteq B'\subseteq A_{gc}$, the reduction in this gap is
\[
 H(\Theta\mid X_B)-H(\Theta\mid X_{B'})
 =I(\Theta;X_{B'\setminus B}\mid X_B)\ge0.
\]
More points can inform the shared parameters. Exact recovery still fails for
finite blocks. Even knowing one cluster center need not determine its parent
group center. When $B$ extends to another cluster, its contributing parameter
family changes, so the displayed monotonicity argument no longer applies as
written. This example motivates block witnesses and nonzero recovery errors;
it does not establish a chosen uniform error bound.

\paragraph{Source.}
The motivating hierarchy is from Gillen and Chiang,
\emph{A summary of
Condensation and its relation to Natural Latents}~\citep{gillen2026summary}, \S1.1.1. The finite law,
illustration, and exact score comparisons above are this course's adaptation.

\section{Condensation in context}
\subsection{What is being characterized?}
Fix jointly distributed observables $X_1,\ldots,X_n$. We organize information
into latents $Y_A$, indexed by nonempty subsets $A$ of the observable indices.
The address specifies which observations the latent is assigned to. Each $X_i$
is determined by its contributing collection $(Y_A:i\in A)$.

LVM alone leaves considerable freedom. The additional reconstruction
and Markov conditions constrain where shared information can reside. The
correspondence theorem compares two models satisfying these conditions: their
collections above the same address determine one another, exactly or with a
controlled conditional-entropy error.

This provides a conditional sense of objectivity. Once the observables and the
conditions are fixed, agreement in information content need not come from a
shared vocabulary. It follows from each model's relationship to the observations.
The theorem does not choose the observables. Changing the questions changes the
problem, and may change the information that must correspond.

Nor is each latent automatically a human concept. A model includes residual
information needed to determine observations. The interpretation of its parts
as concepts is an additional connection between the mathematics and the uses of
a latent model.

\subsection{Connections to prior work}
\label{sec:context-connections}
\subsubsection{Sufficient statistics}
For a fixed joint law of $(X,V)$, a statistic $T=f(X)$ is sufficient for predicting
$V$ when
\[
 I(X;V\mid T)=0.
\]
Knowing $X$ then gives no predictive information about $V$ beyond $T$.
For finite alphabets, grouping values $x$ with the same conditional distribution
$P(V\mid X=x)$ gives a minimal such statistic, on the support of $X$.
This differs from parametric sufficiency, where the conditional law of
the sample given a statistic must be independent of the unknown parameter.
Both formulations ask which distinctions must be retained for an inference task.

Predictive sufficiency does not say that $T$ determines the realized value of
$V$. Condensation's LVM condition does demand determination, using the whole
contributing collection, including residual variables. A sufficient
statistic can itself be a vector or another structured object.

A sufficient statistic is by construction a function of $X$. A latent need not
be, because reconstruction is asked for only up to a defect: that leaves room
for latents which are parameters of a law rather than functions of the
observations. The additional
question here is how to organize a family of latent parts across many specified
observables. See \citet{fisher1922mathematical}
and Shalizi's account of predictive sufficiency~\citep{shalizi2019prediction}.

\subsubsection{Common information}
Two older questions distinguish requirements that appear together in the
perfect case. The maximal deterministic common part of $X_1,X_2$ is a variable
$K$ recoverable from either observation, through which every other common
function factors. Its entropy is the G\'acs--K\"orner common information.
Wyner's common information instead minimizes $I(X_1,X_2;W)$ over extensions
satisfying $X_1\perp X_2\mid W$. Such a $W$ need not be separately recoverable
from the observations.

If $W$ is both separately recoverable and screens off dependence, then
\[
 I(X_1;X_2)=H(W)+I(X_1;X_2\mid W)=H(W).
\]
Moreover, every common function $K$ is determined by $W$: the two-witness lemma,
with witnesses $X_1,X_2$ and common part $W$, gives $H(K\mid W)=0$.
Thus this special case is closely connected to common information in the sense
of G\'acs and K\"orner.

It is restrictive. If $X$ is a fair bit and $V=X\mathbin{\oplus}N$ for independent
$N\sim\operatorname{Bernoulli}(p)$, $0<p<1/2$, all four pairs have positive
probability. Any common function must therefore be constant, although
$I(X;V)=1-h_2(p)>0$, where
$h_2$ is the binary entropy function of \cref{sec:existence}. Recoverable common information and
mutual information differ.

Common/private source coding already studies structured organization
of information. Condensation's particular focus is a family of subset-addressed
latents and correspondence under conditions on that family.
See \citet{gacs1973common}
and Liu, Xu and Chen's extension of Wyner's common information~\citep{liu2010common}.

\subsubsection{Computational mechanics}
Computational mechanics groups past histories according to their conditional
distributions over futures. Its causal states have predictive optimality,
minimality, and uniqueness properties. This is an important precedent for
characterizing such a construction through informational requirements.

The present setup starts from a family of jointly distributed questions rather
than a distinguished past/future split. It characterizes latent collections
associated with different observable subsets.
See \citet{shalizi2001computational}.

\subsubsection{Graphical models and the information bottleneck}
The ordered Markov condition in condensation is a graphical-model condition.
A subset-ordered directed graph places larger addresses above smaller ones.
Under the corresponding factorization, the latent joint law is
\[
 P(y)=\prod_{A\ne\varnothing}P(y_A\mid y_{\supsetneq A}).
\]
A factorization alone does not ensure that observations recover the latents,
or that different latent models correspond. Those are additional requirements.
Graphical language also does not make the result a causal-identification theorem.
Eisenstat discusses this connection explicitly in Definition 5.8 and Example 6.3
of the condensation paper~\citep{eisenstat2025condensation}.

The information bottleneck compresses an input $X$ into $T$ while retaining
information about a specified relevance variable $V$. A standard objective is
\[
 I(X;T)-\beta I(T;V), \qquad T\perp V\mid X.
\]
It provides an established distribution-based approach to useful summaries of
the data.
Our condensation problem instead involves many addressed latents, LVM, and
conditions across observable sets. An optimum of a generic
bottleneck objective does not automatically satisfy these conditions.
See \citet{tishby1999information}.

\subsection{What information theory leaves out}
\subsubsection{An encoding and access to its inverse}
Let $X$ be uniform on $\{1,\ldots,N\}$ and let $Y=\pi(X)$ for a fixed
permutation $\pi$. Then
\[
 H(X\mid Y)=H(Y\mid X)=0.
\]
This is exact correspondence whether $\pi$ is the identity or a large
unstructured lookup table. The entropies do not charge for describing the table,
learning it, storing it, or using it. For a finite table an inverse exists and is
computable; that does not make it available to a particular learner.

For a concrete learning distinction, draw $\pi$ uniformly from all permutations
and reveal $m<N$ distinct input--output pairs. Conditional on those pairs, the
image of a previously unseen input is uniform over the $N-m$ unused outputs. The best success
probability without further information is $1/(N-m)$. For every fixed permutation
the correspondence entropy is zero, yet a learner with few examples cannot
predict an unseen pairing. These are different probability models: one fixes
the permutation; the other describes uncertainty about it.

A separate question concerns computation even when a transformation has a short
known description. Entropic bounds alone give no running-time bound for an
inverse. Establishing computational difficulty requires an additional argument
or assumption; it does not follow from this permutation example.

\subsubsection{Rare events and average error}
Let $E\sim\operatorname{Bernoulli}(p)$. Define $R=0$ when $E=0$ and, when
$E=1$, let $R$ be uniform over $M$ other values. Since $R$ reveals $E$,
\[
 H(R)=h_2(p)+p\log_2 M.
\]
For $p=10^{-6}$ and $M=2^{10^6}$, the rare branch contributes one full bit.
An independent sample of size $m$ misses that branch with probability
$(1-p)^m$. Large informational contributions can be absent from a small
dataset.

Conversely, suppose a variable $T$ records $E$, and records the value of a fair
bit $B$, independent of $E$, only when $E=0$. Then
\[
 H(B\mid T)=p.
\]
The average uncertainty is small, but conditional on the rare event $E=1$, the
bit is entirely unknown. If that event is a critical decision, its importance
is not measured by its probability alone. Entropy is an average informational
quantity, not a measure of the consequences of errors.

These examples distinguish informational equivalence, statistical access, and
performance on important cases. A useful translation may need guarantees about
all three.

\subsection{The connection to AI safety}
\label{sec:context-safety}
\begin{enumerate}
\item \textbf{Corresponding information.} Under the reconstruction and Markov
conditions, the theorem gives correspondence between specified collections in
two models of the same observables.
\item \textbf{Concepts of actual agents.} To apply that result, we need to
identify suitable variables in actual agents and establish the relevant
conditions, including how their observables relate. The theorem does not show
that learning procedures produce such models.
\item \textbf{A usable translation.} Informational correspondence could support
identifying which of an AI's latents carries something humans care about. We
must still find
the correspondence with available data and computation, and connect it to the
particular human distinction of interest.
\item \textbf{Interpretation and specification.} A usable translation could help
interpret predictions or express objectives in an AI's world model. This requires
preserving relevant distinctions when circumstances or models change.
\item \textbf{Safer systems.} Better interpretation and specification could improve
oversight and alignment. Sharing concepts does not entail sharing values,
reporting faithfully, or following an intended objective.
\end{enumerate}
The first step is the mathematical result of this lesson. The subsequent steps
state a research program and the additional work required to connect it to
safety.

\subsection{Discussion}
\begin{enumerate}
\item \textbf{Choosing the questions.} Consider two agents describing the same
physical system, one using local measurements and another using a different
collection of measurements. What would have to be established before applying
a theorem that assumes the same observables? Distinguish changing value labels
from changing which distinctions an observable makes. Identify one restriction
on observable choice that seems justified for a concrete application.
\item \textbf{What would count as a useful translation?} Suppose two latent tuples $U,V$
satisfy $H(U\mid V)=H(V\mid U)=0$. Specify an additional test
that would make this correspondence useful for interpreting a learned model.
State what access the test requires: paired samples, interventions, model
internals, or a known decoder. Which part of the permutation example would
that access resolve?
\item \textbf{From correspondence to oversight.} Choose one human distinction
that matters for evaluating an AI's behavior. Trace the steps from a
correspondence theorem to an oversight procedure that uses that distinction.
Name a failure that could occur despite exact informational correspondence,
and state what additional evidence would address it.
\end{enumerate}

\section{Quiz}
Take the quiz at
\href{https://docs.google.com/forms/d/e/1FAIpQLSdTAs3Dq6dOeoju0TsyfaVcEcRdZaNgpAe_FMC7rCR3fohOSw/viewform}{this
link}. The questions are reproduced here because the form cannot typeset
mathematics.

\subsection{Questions}
Use base-two logarithms. Throughout, $S,P,Q$ are independent fair bits and
$X_1=(S,P)$, $X_2=(S,Q)$, $X_3=(P,Q)$, with latents $Y_{\{1,2\}}=S$,
$Y_{\{1,3\}}=P$ and $Y_{\{2,3\}}=Q$.

\begin{enumerate}
\item What is $I(X_1;X_2)$?
\begin{itemize}
\item $0$
\item $1$
\item $2$
\item $3$
\end{itemize}

\item Same model. What is $H(Y_{\{1,2\}}\mid X_1)$?
\begin{itemize}
\item $0$
\item $1/2$
\item $1$
\item $2$
\end{itemize}

\item Now put $Y_{\{i\}}=X_i$ at each singleton address, with every other slot
constant. Which condition fails?
\begin{itemize}
\item the latent variable model condition (LVM)
\item the reconstruction condition
\item the ordered Markov condition
\item none of them fails
\end{itemize}

\item Complete the two-witness bound
$H(U\mid M)\le H(U\mid M,P)+H(U\mid M,Q)+{}$\,?
\begin{itemize}
\item $I(P;Q)$
\item $I(P;Q\mid M)$
\item $H(P\mid Q)$
\item $0$
\end{itemize}

\item In the correspondence theorem with $|A|=k$, suppose every reconstruction
error and every Markov defect is at most $\varepsilon$. What is the bound on
$H(Y_{\Up A}\mid Z_{\Up A})$?
\begin{itemize}
\item $k\varepsilon$
\item $(k-1)\varepsilon$
\item $(2k-1)\varepsilon$
\item $2k\varepsilon$
\end{itemize}

\item Two corresponding latent collections satisfy $H(U\mid V)=H(V\mid U)=0$.
Which of these does that \emph{not} establish?
\begin{itemize}
\item each collection is a function of the other almost everywhere
\item neither collection retains uncertainty about the other
\item an efficient procedure for translating between them
\item the two collections carry the same information
\end{itemize}
\end{enumerate}

\section{Further reading}
\label{sec:further-reading}
\subsection{Condensation}
\label{sec:reading}
The required mathematics is contained in this sheet and the lectures.
\begin{itemize}
\item Sam Eisenstat, \emph{Condensation: A Theory of Concepts} (2025).
Definitions and exact characterization: Sections 4--5. Correspondence:
Theorems 5.15 and 6.8; two-witness inequality: Lemma 6.4
\citep{eisenstat2025condensation}.
\item Satya Benson, \emph{Relating Almost Perfect Condensation Conditions and
the Conditioned Score}. The score/condition bridge and its reconstruction
coverage qualification
\citep{benson2026relating}.
\item For the wider statistical context: Cover and Thomas,
\emph{Elements of Information Theory}~\citep{cover2006elements}; Lauritzen, \emph{Graphical Models}~\citep{lauritzen1996graphical};
Shalizi and Crutchfield, \emph{Computational Mechanics: Pattern and Prediction,
Structure and Simplicity} (2001)~\citep{shalizi2001computational}.
\end{itemize}

\subsection{Condensation in context}
The sources cited above support different parts
of the discussion. For the mathematical development, start with
Eisenstat's \emph{Condensation: A Theory of Concepts}~\citep{eisenstat2025condensation},
particularly the exact and quantitative correspondence results.
For the score connection, see
\emph{Relating Almost Perfect Condensation Conditions and the Conditioned Score}~\citep{benson2026relating}.
The score decomposition connects excess conditioned code length to dependence
and recoverability defects; the converse requires attention to which observation
sets the reconstruction assumptions cover.

\bibliographystyle{plainnat}
\bibliography{biblo}
\end{document}
